\problemstart
\addcontentsline{toc}{section}{1.6　拘束条件つき運動}

\begin{problemBox}{問題 1.6\quad 拘束条件つき運動\hspace{1.2em}{\small\mdseries 〔標準〕}}
天井に固定された定滑車に1本の軽い糸Aを通し，その一端に質量 $m_1$ の物体1を吊り下げ，他端に質量 $m_{\text{p}}$ の動滑車を吊り下げる。さらに，その動滑車に通したもう1本の軽い糸Bの一端に質量 $m_2$ の物体2を吊り下げ，他端に質量 $m_3$ の物体3を吊り下げる。
\par
鉛直下向きを正の方向とする座標軸を設定し，天井の定滑車の中心を原点 O としたときの，物体1，物体2，物体3，および動滑車の位置座標をそれぞれ $x_1, x_2, x_3, x_{\text{p}}$ とする。糸Aの長さを $L_{\text{A}}$ ，糸Bの長さを $L_{\text{B}}$ とし，滑車の半径および糸の質量は無視できるものとして，以下の問いに答えよ。

\begin{enumerate}
  \item 糸Aの長さが一定であるという条件から，物体1の位置 $x_1$ と動滑車の位置 $x_{\text{p}}$ の間に成り立つ拘束条件式を示せ。また，この式を時間で2回微分することにより，物体1の加速度 $a_1 = \ddot{x}_1$ と動滑車の加速度 $a_{\text{p}} = \ddot{x}_{\text{p}}$ の間に成り立つ関係式を求めよ。

  \item 糸Bの長さが一定であるという条件から，動滑車を基準とした物体2と物体3の相対位置を考慮し， $x_2, x_3, x_{\text{p}}$ の間に成り立つ拘束条件式を示せ。

  \item (1) と (2) の結果から動滑車の座標 $x_{\text{p}}$ およびその微分項を消去し，3つの物体の加速度 $a_1, a_2 = \ddot{x}_2, a_3 = \ddot{x}_3$ の間に成り立つ拘束条件関係式を導出せよ。
\end{enumerate}
\end{problemBox}

\solutionheading
\begin{solutionparts}

\solutionpart{(1)}
原点から物体1までの距離は $x_1$，動滑車までの距離は $x_{\text{p}}$ です。糸Aの長さが $L_{\text{A}}$ で一定であるため，次の拘束条件が成り立ちます。
\begin{equation*}
x_1 + x_{\text{p}} = L_{\text{A}}
\end{equation*}
この式の両辺を時間 $t$ で2回微分します。右辺の糸の長さ $L_{\text{A}}$ は定数であるため，微分すると 0 になります。
\begin{equation*}
\ddot{x}_1 + \ddot{x}_{\text{p}} = 0
\end{equation*}
加速度の記号を用いて整理すると，次の関係式が得られます。
\begin{equation*}
a_1 + a_{\text{p}} = 0
\end{equation*}

\solutionpart{(2)}
糸Bは動滑車に吊り下げられています。動滑車の位置座標 $x_{\text{p}}$ を基準にすると，そこから測った物体2の相対的な位置は $x_2 - x_{\text{p}}$ であり，物体3の相対的な位置は $x_3 - x_{\text{p}}$ と表されます。この2つの長さの合計が，糸Bの全長 $L_{\text{B}}$ に一致しなければなりません。したがって，求める拘束条件式は次のようになります。
\begin{equation*}
(x_2 - x_{\text{p}}) + (x_3 - x_{\text{p}}) = L_{\text{B}}
\end{equation*}
この式の左辺の括弧を外して整理します。
\begin{equation*}
x_2 + x_3 - 2x_{\text{p}} = L_{\text{B}}
\end{equation*}

\solutionpart{(3)}
(2) で得られた拘束条件式の両辺を時間 $t$ で2回微分します。糸の長さ $L_{\text{B}}$ は一定であるため，右辺は 0 になります。
\begin{equation*}
\ddot{x}_2 + \ddot{x}_3 - 2\ddot{x}_{\text{p}} = 0
\end{equation*}
各項を加速度の記号に置き換えます。
\begin{equation*}
a_2 + a_3 - 2a_{\text{p}} = 0
\end{equation*}
ここで，(1) で導出した加速度の関係式 $a_1 + a_{\text{p}} = 0$ から，動滑車の加速度について整理した式 $a_{\text{p}} = -a_1$ を上の等式に代入します。
\begin{equation*}
a_2 + a_3 - 2(-a_1) = 0
\end{equation*}
左辺の符号を展開して整理します。
\begin{equation*}
2a_1 + a_2 + a_3 = 0
\end{equation*}
したがって，3つの物体の加速度は $2a_1+a_2+a_3=0$ を満たします。

\end{solutionparts}
